\pagestyle{fancy}
\fancyhead[l]{\autorUO}
\fancyfoot[l]{\asignaturaAbbr}
\fancyfoot[r]{\fecha}

\section{Sistema de desarrollo} \label{sec:intro_dev_system}
Para trabajar de manera eficiente con \pytorch\ y aprovechar la aceleración mediante GPU, es necesario contar con un sistema de
desarrollo adecuado. Este abarca tanto el sistema operativo como las herramientas necesarias para gestionar dependencias y bibliotecas externas.

\subsection{Sistema operativo} \label{sec:intro_os}
Aunque \pytorch\ ofrece compatibilidad multiplataforma, el ecosistema de desarrollo más utilizado para aplicaciones de computación acelerada
por GPU es Linux. Las distribuciones de Linux proporcionan un entorno muy flexible y personalizable, ampliamente compatible con las bibliotecas
y controladores necesarios para trabajar con \pytorch. A lo largo de este documento se asumirá el uso de un sistema operativo basado en Linux
—concretamente Fedora Linux 42\footnote{\url{www.fedoraproject.org}}—, aunque las instrucciones pueden adaptarse fácilmente a otros sistemas operativos como Windows o macOS.

\subsection{Visual Studio Code} \label{sec:intro_vscode}

\subsection{Jupyter Notebook} \label{sec:intro_jupyter}